%%%%%%%%%%%%%%%%%%%%%%%%%%%%%%%%%%%%%%%%%%%%%%%%%%%%%%%%%%%%%%%%%%%%%%%%%%%%%%%%%%%%%%%%%
%
% Modelo para poster da Semana de Estudos de Física e Matemática da UEMASUL
% altere de acordo com suas informações
%
% criado por: Clovis Caface
% portado para o Overleaf por Clovis Caface
%
% license by LaTeX Project Public License 1.3c
%
% LEIA COM ATENCAO
%
% Este projeto está composto por texto gerado automaticamente nos comandos \blindtext[*],
% portanto retire o comando e no lugar coloque o texto de sua apresentação.
%
% Lembre-se que o poster deve estar em medida de 120x90cm 
%%%%%%%%%%%%%%%%%%%%%%%%%%%%%%%%%%%%%%%%%%%%%%%%%%%%%%%%%%%%%%%%%%%%%%%%%%%%%%%%%%%%%%%%%

\documentclass{poster_sefim}

\title{Titulo do Projeto}
\author{Primeiro Autor; Segundo Autor; Terceiro Autor; Quarto Autor; Quinto Autor.}
\supervisor{Orientador}
\institute{Universidade Estadual da Região Tocantina do Maranhão}
\date{30 de agosto de 2019}

\logoinstitute{\includegraphics[height=1.2cm, keepaspectratio]{logo/universidade.png} % alterar para a logo de sua instituicao
    % \qquad\includegraphics[height=1.2cm, keepaspectratio]{logo/uemasul.png}
    % retirar o comentário da linha acima caso queira inserir um segundo logo
} % Logo da Instituição 

\begin{document}

    %%%%% Não alterar esta quatro linhas abaixo %%%%%%%
    \begin{frame}{} 
        \titlebox
        \vspace{0.7cm}
    \begin{columns} 
    
    
    %%%%%% inicio coluna 1 %%%%%%%%
    \begin{columnbox} % divisão em coluna
        \begin{block}{\Large Introdução}
            \justifying
            \blindtext[3] % comando para preenchimento de texto gerado automaticamente
        \end{block}
        
        \blindtext[1]
        
        \begin{figure}
            \centering
            \includegraphics[scale=0.1]{logo/bottom.png}
            \caption{Simbolo SEFiM}
            \label{fig:my_label}
        \end{figure}
        
        \begin{block}{\Large Revisão Bibliográfica}
            \begin{table}[]
                \centering
                \begin{tabular}{c|c}
                     lorem & ipsun \\
                     dolor & sit
                \end{tabular}
                \caption{Tabela de Lorem Ipsun}
                \label{tab:my_label}
            \end{table}
        \end{block}
        
    \end{columnbox} % finalização de divisão em coluna
    
     %%%%%% inicio coluna 2 %%%%%%%%
    \begin{columnbox} % divisão em coluna
        \begin{block}{\Large Resultados}
            \justifying
            \blindtext[4] % comando para preenchimento de texto gerado automaticamente
        \end{block} 
        \blindtext[2]
    \end{columnbox} % finalização em coluna
    
     %%%%%% inicio coluna 3 %%%%%%%%
    \begin{columnbox} % divisão em coluna
        \begin{block}{\Large Conclusão}
            \justifying
            \blindtext[3] % comando para preenchimento de texto gerado automaticamente
        \end{block}
        \blindtext[1]
        \begin{block}{\Large Referências}
            \begin{thebibliography}{9}
            % altere os dados abaixo para as suas referencias
                \bibitem{latexcompanion} 
                Michel Goossens, Frank Mittelbach, and Alexander Samarin. 
                \textit{The \LaTeX\ Companion}. 
                Addison-Wesley, Reading, Massachusetts, 1993.
                 
                \bibitem{einstein} 
                Albert Einstein. 
                \textit{Zur Elektrodynamik bewegter K{\"o}rper}. (German) 
                [\textit{On the electrodynamics of moving bodies}]. 
                Annalen der Physik, 322(10):891–921, 1905.
                 
                \bibitem{knuthwebsite} 
                Knuth: Computers and Typesetting,
                \\\textit{http://www-cs-faculty.stanford.edu/\~{}uno/abcde.html}
            \end{thebibliography}
        \end{block}
        
        \begin{block}{\Large Agradecimentos}
            \justifying
            % altere os dados abaixo para os seus agradecimentos
            \begin{itemize}
                \item Primeiro item 
                \item Segundo item
                \item Terceiro item
                \item Quarto item 
            \end{itemize}
        \end{block}
        
    \end{columnbox}

    \end{columns} % finalização em coluna
    \end{frame}
\end{document}

